\documentclass{article}
\usepackage{graphicx}
\usepackage{xeCJK}
\usepackage{fontspec}
\usepackage{amsmath}

\title{ML part}
\author{saiisagod}
\date{2026年1月31日}

\begin{document}

\maketitle

\noindent\textbf{说在前面：}
本文档受限于作者的经验与能力，暂时无法对常考不常考做出可靠清晰的分类。因此建议各位同学以高频基础（自己觉得）为主，结合自身目标岗位与面试反馈取舍补充。以及，不要想着把文档中的内容背下来，用自己的语言能说出来才是最好。以以及，文档现在肯定是不全的，不过我还没开始准备暑期实习，三四月份该文档才会有最终形态。答案大部分都是gpt写的，他答题能力远超我。

\section{欠拟合、过拟合}
\begin{enumerate}
\item \textbf{什么是欠拟合？什么是过拟合？如何解决？}\\
欠拟合指模型过于简单，无法捕捉数据的真实模式，在训练集和测试集上表现都不好；过拟合指模型过于复杂，过度记忆训练集噪声，在训练集表现很好但测试集性能很差。解决欠拟合可以：增加模型复杂度（更复杂模型、增加特征）、减少正则强度、训练更久。解决过拟合可以：增加数据（数据增强）、使用正则化（L1/L2）、Dropout、简化模型结构、使用交叉验证选择合适的超参数、早停等。

\item \textbf{欠拟合会有什么表现？过拟合会有什么表现？}\\
欠拟合：模型学不出规律，训练集表现就差（loss高/准确率低），测试集也差，二者都不好、差距不大。过拟合：模型把训练集“记住了”，训练集表现很好（loss低/准确率高），但测试集明显变差，训练与测试差距很大。

\item \textbf{什么是偏差-方差权衡？}\\
学习算法的期望泛化误差可分解为：误差 = 偏差$^2$ + 方差 + 不可约误差。偏差反映模型表达能力是否足够（偏差大对应欠拟合），方差反映模型对训练数据扰动的敏感程度（方差大对应过拟合）。提高模型复杂度通常会降低偏差、增加方差，因此需要在偏差和方差之间取得平衡，通过模型选择和正则化实现。

\item \textbf{如何判断模型是否存在高方差？}\\
典型现象是：训练集误差很低，但验证集/测试集误差明显偏高，两者差距较大。使用学习曲线（训练集大小 vs 误差）也能判断：若训练误差很低但测试误差不随数据量增加明显下降，则为高方差。应对方式包括：减小模型复杂度、增加正则化、增加训练数据、使用集成方法等。

\item \textbf{为什么深度模型更容易过拟合？}\\
深度模型参数量巨大，函数空间极其丰富，有能力拟合训练数据中的噪声和偶然结构；如果数据量相对不足、正则化不充分，就容易出现训练误差很小而测试误差较大的情况。此外，深网对输入扰动较敏感（如对抗样本），也体现了高方差特性。
\end{enumerate}

\section{L1/L2}
\begin{enumerate}
\item \textbf{L1 正则与 L2 正则的区别？}\\
L1 正则是在损失中加入参数绝对值之和 $\lambda \sum |w_i|$，L2 正则加入参数平方和 $\lambda \sum w_i^2$。L1 会产生稀疏解（部分参数等于 0）（稀疏解这个部分可以去找个图片看看很直观），有特征选择效果；L2 更偏向让所有参数都变小但不为 0，抑制大参数，有平滑效果。几何上，L1 的约束是菱形，L2 是圆形，最优解与坐标轴交点概率不同导致稀疏性差异。\\
关于这个lasso稀疏解，我也是最近学到可以用贝叶斯解释，系数beta的prior选laplace分布。

\item \textbf{Lasso 和 Ridge 的核心差异？}\\
Lasso 回归使用 L1 正则，目标函数是 $\|y - Xw\|_2^2 + \lambda \|w\|_1$；Ridge 使用 L2 正则，目标为 $\|y - Xw\|_2^2 + \lambda \|w\|_2^2$。Lasso 可以将部分系数压到 0，实现自动特征选择；Ridge 更适合特征较多且存在多重共线性时的稳定估计（更稳定更平滑，对共线性更好）。

\item \textbf{为什么 L1 会产生稀疏解？}\\
从几何角度看，L1 约束 $\sum|w_i|\le c$ 是一个尖角很多的多面体（二维是菱形），与同心圆的损失等高线相切时，最容易在坐标轴上发生接触（某些维度为 0）。从优化角度看，L1 正则的梯度在 0 点存在“不可导但有次梯度”的尖点，优化时更倾向于让小的权重直接收缩到 0，从而得到稀疏解。

\item \textbf{什么是特征共线性？怎么解决？}\\
特征共线性指多个特征之间存在强线性相关，导致设计矩阵列之间近线性依赖。这样会使线性模型参数估计不稳定、方差增大、系数解释困难。解决方法包括：删除高度相关的特征；使用 PCA 等降维；引入正则化（如 Ridge）；或使用对多重共线性不敏感的模型（如树模型）。

\item \textbf{线性模型为什么可解释性强？}\\
线性模型假设输出是输入特征的线性加权和，每个特征对应一个系数，可以直观解释为“该特征单位变化对输出的边际影响”，同时没有复杂的非线性叠加结构，决策边界也是线性的。在适当标准化下，系数大小、符号都具有明确含义，因此易于解释。
\end{enumerate}

\section{模型选择与超参}
\begin{enumerate}
\item \textbf{交叉验证的原理是什么？}\\
交叉验证通过将数据集划分为多个互斥子集，轮流将其中一部分作为验证集、其余作为训练集，计算模型在不同划分上的平均性能，以此估计模型在未见数据上的泛化能力。这样可以更充分利用有限数据，减少单一划分带来的偶然性，提高模型选择和超参数调节的稳定性。（一般是5 or 10）

\item \textbf{为什么要用 K 折交叉验证？}\\
K 折交叉验证（如 $K=5$ 或 $10$）在偏差和方差之间取得折中：$K$ 太小（如留出法）方差较大，结果受一次划分影响；$K$ 太大（如留一法）计算开销大且方差也不一定最小。K 折可以重复利用大部分数据训练，每次使用不同子集验证，使评估结果更稳健。

\item \textbf{什么是超参数？什么是模型参数？}\\
模型参数是通过训练数据学习得到的。超参数是在训练前人为设定，通常通过交叉验证等模型选择方法来确定。
\end{enumerate}

\section{深度学习训练技巧（Dropout/归一化/梯度稳定性）}
\begin{enumerate}
\item \textbf{Dropout 的原理是什么？}\\
Dropout 在训练阶段以一定概率随机将神经元输出置零，相当于每次迭代采样一个子网络进行训练。这样可以打破神经元之间的共适应关系，不再依赖某一批神经元。

\item \textbf{Batch Normalization 的作用与原理？}\\
BN 在每个 mini-batch 内对中间层的激活进行标准化（减均值、除以标准差），然后通过可学习的缩放和平移恢复表达能力，可缓解内部协变量偏移。

\item \textbf{LayerNorm、BatchNorm、InstanceNorm 的区别？}\\
BN 在 batch 维度上做归一化，对同一通道在一个批次的样本上统计均值方差；LayerNorm 对单个样本的特征维度做归一化，常用在 RNN/Transformer 中，对 batch 大小不敏感；InstanceNorm 对每个样本、每个通道单独计算均值方差，常用于风格迁移等。不同归一化方式适用于不同结构和任务。

\item \textbf{什么是梯度消失？什么是梯度爆炸？}\\
在深层网络中，误差通过链式法则反向传播，如果梯度项绝对值小于 1，层层相乘会趋近于 0，导致靠前层几乎无法更新，称为梯度消失；若梯度项绝对值大于 1，层数多时梯度会指数级放大，导致参数更新不稳定甚至溢出，称为梯度爆炸。常见解决方法有合适的激活函数（ReLU 系列）、初始化策略、残差结构、梯度裁剪。（关于这些可以看看各种RELU 残差连接（面试官问过我transformer是pre还是post） 梯度裁剪之类的方式）

\item \textbf{什么是梯度裁剪？为什么需要？}\\
梯度裁剪指在反向传播时，对梯度的范数进行上界限制，如若 $\|\nabla\| > \tau$，则将梯度缩放为 $\tau \cdot \nabla/\|\nabla\|$。在 RNN、深层网络中，梯度爆炸会导致参数剧烈震荡甚至数值溢出，梯度裁剪可以稳定训练过程。
\end{enumerate}

\section{激活函数与输出层（ReLU族/Sigmoid/Softmax）}
\begin{enumerate}
\item \textbf{为什么 ReLU 能缓解梯度消失？}\\
sigmoid/tanh 在较大正负输入时梯度趋近于 0，而 ReLU 在正半轴上导数为 1，不会出现饱和区梯度趋近 0 的问题，因此在多层网络中梯度可以较顺畅地向前传播，减轻梯度消失。但 ReLU 也有“死亡 ReLU”问题，因此有 LeakyReLU、PReLU 等变体。

\item \textbf{ReLU、LeakyReLU、ELU、GELU 区别？}\\
ReLU 为 $f(x)=\max(0,x)$，简单高效；LeakyReLU 在负半轴保留一小段斜率 $ax$（$a$ 很小），缓解神经元死亡；ELU 在负半轴采用平滑的指数形式，有更好的平滑性和归一化特性；GELU 是高斯误差线性单元，输出为 $x \cdot \Phi(x)$（$\Phi$ 为标准正态 CDF），常用于 Transformer 等模型，被认为在某些任务上表现在优。

\item \textbf{Sigmoid 的缺点是什么？}\\
Sigmoid 在正负大数值区域梯度趋近于 0，易导致梯度消失；输出不是零均值（范围在 $(0,1)$），会造成后续层激活偏移，影响收敛；Sigmoid 中包含指数运算，计算相对 ReLU 等更慢。因此在隐藏层中已较少使用，只在输出层（如二分类）中仍然常用。

\item \textbf{softmax 为什么要减去最大值？}\\
softmax 为 $\text{softmax}_i(z)=\frac{e^{z_i}}{\sum_j e^{z_j}}$，若 $z_i$ 绝对值较大，$e^{z_i}$ 容易溢出。由于 softmax 对所有输入同时加减同一常数不变，因此通常用 $z_i' = z_i - \max_j z_j$，使最大值为 0，其余为负数，从而保证指数在数值范围内，提升数值稳定性。

\item \textbf{softmax 的梯度形式？}\\
设 $p_i = \text{softmax}_i(z)$，则对 $z_k$ 的偏导为
\[
\frac{\partial p_i}{\partial z_k} =
\begin{cases}
p_i(1 - p_i), & i = k,\\
-p_i p_k, & i \ne k.
\end{cases}
\]
若结合交叉熵损失 $L = -\sum_i y_i \log p_i$，则 $\frac{\partial L}{\partial z_k} = p_k - y_k$，形式非常简洁，便于反向传播实现。
\end{enumerate}

\section{损失函数}
\begin{enumerate}
\item \textbf{什么是交叉熵？为什么常用于分类？}\\
交叉熵衡量两个分布之间的差异，离散情形为 $H(p,q)=-\sum_i p_i \log q_i$。在分类中，$p$ 是真实标签（通常为 one-hot），$q$ 是模型预测概率。最小化交叉熵等价于最大化真实标签的对数似然，也等价于最小化 KL 散度，具有良好的统计性质，并且梯度形式简单（$p-q$），在优化上非常方便。

\item \textbf{为什么分类不用 MSE？}\\
MSE 本质是高斯噪声假设下的最优损失，而分类问题更适合伯努利/多项式分布假设，对应交叉熵。若用 MSE + sigmoid，梯度会包含额外的 Sigmoid 导数项，容易饱和；交叉熵与 sigmoid/softmax 结合时梯度形式为 $p-y$，数值更稳定，收敛更快。

\item \textbf{损失函数有哪些？}\\
回归中常用：MSE、MAE、Huber 损失等；分类中常用：交叉熵、Focal Loss 等；排序中：pairwise/ listwise 损失；生成中：对抗损失、重构损失（L1/L2）、感知损失等。损失函数编码了任务的目标与模型输出的偏好。

\item \textbf{Focal Loss 的作用是什么？}\\
Focal Loss 针对正负样本极度不平衡且易样本过多的问题，通过在交叉熵上乘以 $(1-p_t)^\gamma$，降低对易分类样本的损失权重，放大对难分类样本的关注。这样模型训练时不会被大量简单样本主导，常用于检测中的正负样本极不平衡场景。

\item \textbf{不平衡样本如何处理？}\\
常用方法包括：在数据层面，欠采样多数类、过采样少数类、生成合成样本；在损失层面，使用类别权重或 Focal Loss；在模型层面，调整决策阈值、使用专门针对不平衡问题的算法；在评估层面，使用 AUC、F1、AUPR 等适合不平衡的指标而非简单精度。
\end{enumerate}

\section{经典浅层模型（逻辑回归/SVM/核技巧）}
\begin{enumerate}
\item \textbf{什么是逻辑回归？它能处理非线性吗？}\\
逻辑回归是一种用于二分类的广义线性模型，假设 $P(y=1|x)=\sigma(w^\top x+b)$，其中 $\sigma$ 为 sigmoid。。

\item \textbf{逻辑回归的似然函数是什么（这个不太需要看）}\\
对于样本 $(x_i,y_i)$，$y_i\in\{0,1\}$，$P(y_i=1|x_i)=\sigma(w^\top x_i)$，则总体似然为
\[
L(w)=\prod_i \sigma(w^\top x_i)^{y_i} \bigl(1-\sigma(w^\top x_i)\bigr)^{1-y_i},
\]
log-likelihood 为
\[
\ell(w)=\sum_i \left[y_i \log \sigma(w^\top x_i) + (1-y_i)\log\bigl(1-\sigma(w^\top x_i)\bigr)\right],
\]
通常最大化 log-likelihood 或最小化其负数（即交叉熵损失）。

\item \textbf{为什么逻辑回归要用最大似然？}\\
逻辑回归来自一个明确的概率模型，在这种框架下，最大似然估计是自然的参数估计方法，具有一致性、渐近无偏、渐近正态以及在一定条件下的渐近有效性等良好性质。此外，最大似然对应的损失是凸的，便于优化。

\item \textbf{SVM 的硬间隔和软间隔是什么？}\\
硬间隔 SVM 假设数据完全线性可分，要求所有样本满足 $y_i(w^\top x_i + b)\ge 1$，最大化几何间隔；软间隔 SVM 引入松弛变量 $\xi_i$，允许少量样本违反间隔约束，通过惩罚项 $C\sum \xi_i$ 权衡间隔最大化和误分类，适用于带噪声的实际数据。

\item \textbf{核函数的作用是什么？}\\
核函数 $K(x,z)$ 等价于在高维（甚至无穷维）特征空间中计算内积 $\phi(x)^\top\phi(z)$，这样可以在不显式计算高维映射 $\phi$ 的前提下，使用线性算法（如 SVM）得到非线性决策边界。核技巧的核心是：只要算法只依赖于数据点的内积，就可以用核函数替代。

\item \textbf{常见的核函数有哪些？}\\
常见核函数包括：线性核 $K(x,z)=x^\top z$；多项式核 $K(x,z)=(x^\top z + c)^d$；高斯 RBF 核 $K(x,z)=\exp(-\|x-z\|^2/(2\sigma^2))$；Sigmoid 核 $K(x,z)=\tanh(\alpha x^\top z+\beta)$ 等。不同核对应不同的隐式特征空间和光滑性假设。
\end{enumerate}

\section{树模型与集成学习}
\begin{enumerate}
\item \textbf{关于这章b站上有很多视频，宜配合视频服用}\\
\item \textbf{熵是什么？基尼系数又是什么？}\\
熵（Entropy）是信息论中的一个重要概念，用于衡量系统的不确定性。在机器学习中，熵常用于衡量数据集的纯度。熵的数学表达式如下：
$$
H(S) = -\sum_{i=1}^{n} p_i \log_2 p_i
$$
其中，$p_i$ 表示第 $i$ 类在数据集 $S$ 中出现的概率。熵越大，表示数据集越混杂，纯度越低；熵越小，表示数据集越纯净。

基尼系数（Gini Index）也是衡量数据集纯度的指标，常用于决策树算法。基尼系数的公式如下：
$$
Gini(S) = 1 - \sum_{i=1}^{n} p_i^2
$$
其中，$p_i$ 表示第 $i$ 类的概率。基尼系数越小，数据集越纯净。
\item \textbf{信息增益是什么？（用于分类树）}\\
信息增益（Information Gain）是用于衡量某个特征对数据集分类能力提升的指标，常用于决策树的特征选择。其计算方法是选择某个特征进行划分后，数据集的熵减少了多少。信息增益的公式如下：
$$
\text{Information Gain}(S, A) = H(S) - \sum_{v \in \text{Values}(A)} \frac{|S_v|}{|S|} H(S_v)
$$
其中，$H(S)$ 表示划分前数据集 $S$ 的熵，$A$ 是特征，$\text{Values}(A)$ 表示特征 $A$ 的所有可能取值，$S_v$ 表示在特征 $A$ 取值为 $v$ 时的数据子集，$H(S_v)$ 是该子集的熵。信息增益越大，说明该特征对分类的提升越显著。
\item \textbf{终止条件可以是什么？树的过拟合是什么？什么是剪枝？}\\
常见的终止条件包括（Pre-pruning）：1、树的最大深度2、没有可以使用的feature3、信息增益小于threshold4、节点samples的数量小于总样本数量*0.05（注意：终止条件通常就是预剪枝的具体表现形式Pre-pruning）。树的过拟合（Overfitting）是指决策树模型过度学习训练数据中的“噪声”，导致模型非常复杂，对训练数据拟合很好，但泛化能力差，对新数据表现不佳。剪枝是处理过拟合的一种方法，包括有pre pruning以及 post pruning，后剪枝：1、 减少误差的后剪枝方法(reduced-error post-pruning): 把数据集分为训练集和验证集，尝试剪掉每一个节点并测试在验证集上的结果。每次都贪心的剪掉那个带来最大验证集上准确率提升的节点。（Split data into training set and validation set, do until further pruning is harmful: - Evaluate impact on validation set of pruning each possible node - Greedily remove the one that most improves validation set accuracy mostly.）2、规则后剪枝（常用于C4.5）步骤如下：

1. 把决策树转化为一系列规则，每个规则表示从根到叶节点的一条路径。例如：
   \texttt{if (outlook=sunny) $\wedge$ (humidity=high) then playTennis = no}

2. 对每条规则，逐步去除规则中的前提条件（即从规则中移除某些节点），判断去除条件后规则在验证集上的准确率是否提升。如果提升，则移除该条件。例如：
   \texttt{if (outlook=sunny) then playTennis = no}

3. 按照规则的估计准确率（estimated accuracy）对所有规则进行排序。

4. 在决策树预测时，使用所有规则，在样本满足多个规则时，选择准确率最高的那条规则进行预测。

这种方法的目的是：通过简化规则，去掉不必要的条件，让决策树变得更简单，减少过拟合，提高在新数据上的表现。
\item \textbf{ID3有什么缺点}\\
无法处理连续值，没有预处理缺失值，没有剪枝。

\item \textbf{C4.5}\\
C4.5决策树算法中，信息增益比（Gain Ratio）是用来选择最佳划分特征的指标。它是在ID3算法的信息增益基础上进行改进，解决了信息增益偏向于取值较多特征的问题。

信息增益比的计算方法如下：

1. 先计算特征 $A$ 的信息增益 $IG(S, A)$：
   $$
   IG(S, A) = H(S) - \sum_{v \in \text{Values}(A)} \frac{|S_v|}{|S|} H(S_v)
   $$
   其中 $H(S)$ 是数据集 $S$ 的熵，$S_v$ 是特征 $A$ 取值为 $v$ 时的数据子集。

2. 再计算特征 $A$ 的分裂信息（Split Information）：
   $$
   SI(S, A) = -\sum_{v \in \text{Values}(A)} \frac{|S_v|}{|S|} \log_2 \frac{|S_v|}{|S|}
   $$

3. 最后，信息增益比公式为：
   $$
   \text{Gain Ratio}(S, A) = \frac{IG(S, A)}{SI(S, A)}
   $$

信息增益比能够避免信息增益偏向于取值多的特征，使特征选择更加合理。但是这样的话不不就偏向取值少的特征了吗，所以这里加了一条：但信息增益比对取值较少的特征有所偏好（分母越小，整体越大），因此 C4.5 并不是直接用信息增益比最大的特征进行划分，而是使用一个启发式方法：先从候选划分特征中找到信息增益高于平均值的特征，再从中选择信息增益比最高的。
（既要增益高又要增益比高）

2、后剪枝使用rule post pruning

3、增加处理连续特征的方法 -- 连续特征离散化
需要处理的样本或样本子集按照连续变量的大小从小到大进行排序
假设该属性对应的不同的属性值一共有N个,那么总共有N−1个可能的候选分割阈值点,每个候选的分割阈值点的值为上述排序后的属性值中两两前后连续元素的中点,根据这个分割点把原来连续的属性分成两类

4、缺失值的处理

特征缺失时的划分  
   当某个样本在某特征上的取值缺失时，C4.5不会简单地丢弃该样本，而是采用“概率分配”的方式。具体做法是：
   - 根据该特征已知取值的样本分布，计算每个分支的概率。
   - 将缺失特征的样本按照这些概率，分配到各个分支，并在后续统计时采用加权方式。

缺失值对信息增益计算的影响  
   在计算信息增益或信息增益比时，C4.5会：
   - 只用特征值已知的样本来计算信息增益。
   - 计算时，对缺失特征的样本进行加权处理，使其对总信息增益的影响按概率分配。

\item \textbf{CART classification and regression}\\

用基尼系数，同时判断只有yes or no 

具体步骤为：

   - 对于每个多分类特征，CART会尝试将所有类别分成两组（例如A和B/C，A/B和C等），枚举所有组合。
   - 计算每种组合的基尼系数，选择最优的二叉划分。
   - 重复上述过程，直到满足终止条件。

连续值同理

优势是：

CART使用二叉树来代替C4.5的多叉树，提高了生成决策树效率
C4.5只能用于分类，CART树可用于分类和回归(Classification And Regression Tree)
CART 使用 Gini 系数作为变量的不纯度量，减少了大量的对数运算
ID3 和 C4.5 只使用一次特征，CART 可多次重复使用特征

\item \textbf{GBDT 和随机森林的区别？}\\
随机森林是基于 Bagging 的集成，每棵树独立训练，通过对样本和特征的随机子采样来降低方差，最终对各树输出做平均/投票。GBDT 是基于 Boosting 的集成，树是逐棵按顺序训练的，每一棵树拟合前一轮残差（或梯度），从而不断降低偏差。前者偏向降低方差、稳定性高；后者偏向降低偏差、表达能力更强但更易过拟合。

\item \textbf{什么是 Boosting？什么是 Bagging？}\\
Bagging（Bootstrap Aggregating）通过对训练集进行有放回采样，训练多个基学习器，再对它们输出平均/投票，以降低方差。Boosting 则是序列式集成，后一个学习器重点学习前一个学习器做错的样本（如 AdaBoost）或残差（如 GBDT），逐步提升整体性能，更偏向降低偏差。

\item \textbf{AdaBoost 的核心思想是什么？}\\
AdaBoost 为每个训练样本维护一个权重，每一轮训练一个弱分类器（分类器是根据错误率最低的分割方式来生成），根据上一轮分类错误情况更新样本权重：被错分的样本权重升高，使得后续分类器更关注难样本。最终通过对弱分类器加权投票形成强分类器。从理论上，AdaBoost 近似在最小化指数损失。

\item \textbf{树模型如何处理缺失值（注意只是label缺失不是特征缺失）？}\\
常见做法包括：在划分时将缺失值样本同时考虑左右子树的增益，选择最优分裂；或者在训练时学习一个“默认方向”，当样本在该特征上缺失时，按照历史统计信息决定进入左子树或右子树（如 XGBoost 中的缺失方向）。也可以在建树前做缺失值填补，但现代树模型一般内置更合理的处理机制。

\item \textbf{决策树如何选择划分点？}\\
对于分类问题，常用的信息增益、信息增益率或基尼指数；对于回归问题，常用最小化均方误差（MSE）或方差。算法会遍历每个特征及其候选分割点，计算划分后子节点的加权指标（如信息熵）是否下降最多，从而选择最优划分。

\item \textbf{信息增益、信息增益率、基尼指数区别？}\\
信息增益是划分前后熵的减少量，偏向取值数目较多的特征；信息增益率对信息增益做归一化，除以特征熵，缓解多值偏好（如 C4.5）。基尼指数衡量分类的不纯度（$1-\sum p_k^2$），与熵类似但计算更简单，CART 常用。三者的选择更多是算法约定和工程习惯。

\item \textbf{为什么深度树更容易过拟合？}\\
树深度越大，叶子节点越多，模型复杂度越高，可以在训练集上拟合非常细致的划分，甚至把单个样本分到独立叶子，导致记住噪声、泛化能力下降。深树对训练数据的小扰动也更敏感，方差大，因此常通过限制最大深度、最小叶子样本数、剪枝等方式控制复杂度。

\item \textbf{什么是剪枝？有哪些剪枝方法？}\\
剪枝是降低决策树复杂度的过程，避免过拟合。预剪枝是在建树过程中设置停止条件，如限制最大深度、最小样本数、不纯度降低不足等，就停止继续划分。后剪枝是先生成一棵较大的树，再自底向上尝试合并某些子树，若合并后验证集误差不升反降，则保留合并，从而得到更简洁的树。

\item \textbf{XGBoost 中的二阶导数用于什么？}\\
XGBoost 使用二阶梯度提升框架，将损失对当前模型输出的泰勒展开到二阶，用一阶导数 $g_i$（梯度）和二阶导数 $h_i$（Hessian）共同构成目标函数的近似。划分时评价每个节点的增益时，会用到 $G=\sum g_i$ 和 $H=\sum h_i$，从而更精确地估计不同划分对损失的影响，使得优化更高效、更稳定。
\item \textbf{RF和GBDT有什么区别？}\\

1、集成学习：RF属于bagging，GBDT是boosting，一个是并行一个是串行，bagging优化的是方差，boosting优化的是偏差，RF不容易过拟合，boosting容易过拟合。

2、训练样本：RF是有放回抽样，GBDT训练使用全部样本

3、数据敏感性：RF对异常值敏感，GBDT对异常值不敏感

\item \textbf{LR和GBDT有什么区别？什么时候LR好？}\\
LR是线性模型，可解释性强，能并行。GBDT非线性，特征表达能力强，但是容易过拟合。高维稀疏场景下LR好，带正则项的LR不容易过拟合。

\item \textbf{XGBOOST和GBDT有什么区别？}\\
1、XGBOOST对叶子节点和深度加入了惩罚项
2、xg是二阶，GBDT是一阶
3、XG支持线性分类器
4、可以自己GPT一下
5、GBDT前剪枝，XGBOOST是后剪枝，先生成完整树再回溯剪枝，效果更优。
\item \textbf{XGBoost 和 LightGBM 有什么区别？}\\
LightGBM 在 XGBoost 的基础上进一步优化：采用基于直方图的高效算法；使用 GOSS（按梯度大小采样）；使用基于叶子（leaf-wise）的生长策略并限制最大深度；支持高效处理大规模类别型特征。整体上 LightGBM 在大规模、高维场景下通常比 XGBoost 更快但更易过拟合，需要更仔细的正则化。
\item \textbf{XGBoost行采样列采样}\\
\item \textbf{LightGBM的GOSS优化}\\
\item \textbf{LightGBM 的直方图加速？}\\
LightGBM 首先将连续特征离散到有限个桶，统计每个桶中样本的梯度和 Hessian 累积值。划分节点时只需要在离散的桶边界上枚举，以及使用整数计数代替浮点运算。

\end{enumerate}

\section{无监督学习K-means}
\begin{enumerate}
\item \textbf{K-means 聚类的目标函数是什么？}\\
K-means 试图最小化簇内平方误差：
\[
J = \sum_{k=1}^K \sum_{x_i \in C_k} \|x_i - \mu_k\|^2,
\]
其中 $C_k$ 为第 $k$ 个簇，$\mu_k$ 为该簇的质心。算法通过迭代“分配样本到最近质心”和“更新质心为簇内均值”两个步骤，直到收敛。

\item \textbf{为什么 K-means 需要多次初始化？}\\
K-means 是一个非凸优化问题，标准算法只能保证收敛到局部最优解，初始质心不同会得到不同的聚类结果。为了找到更好的局部最优，通常进行多次随机初始化，或使用 k-means++ 这样的初始化策略优先选择彼此距离较远的初始中心。

\item \textbf{K-means 容易出现什么问题？}\\
主要问题包括：对初始值敏感；只能发现“球状”簇，对非凸簇形状表现不好；对于簇大小差异很大的情况效果较差；对异常点敏感（异常点可能被作为质心）；需要事先指定簇数 $K$。在高维空间中，由于距离度量不可靠，也会影响效果。
\end{enumerate}

\section{优化与训练策略（优化器/warmup/batch）}
\begin{enumerate}
\item \textbf{Adam 与 SGD 的区别？}\\
SGD 使用全局统一学习率，对每个参数更新大小相同；Adam 为每个参数维护一阶矩（动量）和二阶矩（梯度平方的指数滑动平均），从而根据历史梯度自适应调整每个参数的学习率。Adam 在训练初期收敛更快、对学习率不敏感，但在某些任务上泛化可能略逊于精心调参的 SGD。

\item \textbf{Adam 为什么比 SGD 更快？}\\
Adam 通过动量累积方向、二阶矩调整步长，能在陡峭方向减小步幅、在平坦方向加快前进，从而在复杂非凸损失面上更高效地寻找下降方向。初期阶段尤其明显，常常无需仔细调参就能快速降低训练损失。

\item \textbf{Adam 为什么可能泛化差？}\\
自适应优化器使参数更新步长不均匀，可能导致模型偏向某些“尖锐极小值”，而这类解的泛化性能往往不如 SGD 易到达的“平坦极小值”。此外，Adam 的自适应机制在长期训练中可能导致有效学习率下降过快，不利于探索更好的参数空间。实际中常用 AdamW、逐步切换到 SGD 等策略。

\item \textbf{什么是 warmup？}\\
Warmup 指在训练初期从较小的学习率开始，逐步增加到目标学习率，缓解一开始使用大学习率导致的不稳定（尤其是在深层网络、BatchNorm、Transformer 等中）。Warmup 可以看作一种“热身”，让网络参数先进入一个合理区域，再进行大步更新。

\item \textbf{什么是 mini-batch？大小如何影响训练？}\\
Mini-batch 是每次迭代用于计算梯度的一小部分样本。Batch 越大，梯度估计越接近真实梯度、更新更稳定，但内存占用高、单次迭代慢；Batch 越小，梯度噪声大，但有助于跳出局部极小值、提高泛化。实际中常在硬件允许的范围内选择尽量大的 batch，并配合学习率调整。

\item \textbf{常见优化器包括哪些？}\\
常见优化器包括：标准 SGD，带动量的 SGD（Momentum）、Nesterov 动量；自适应学习率算法如 AdaGrad、RMSProp、Adam、AdamW、AdaDelta 等。不同优化器在收敛速度、对学习率的敏感度和泛化性能上各有特点。
\end{enumerate}

\section{DL结构}
\begin{enumerate}
\item \textbf{padding、stride、dilation 的影响？}\\
Padding 控制边界是否补零，影响输出尺寸和边界信息保留；Stride 控制滑动步长，步长越大，输出越小，相当于下采样；Dilation 控制卷积核元素之间的间隔，增大感受野而不增加参数量，常用于语义分割等需要大感受野的任务。

\item \textbf{为什么要使用残差结构？}\\
残差结构引入 identity shortcut，使网络学习残差映射 $F(x)=H(x)-x$ 而非直接学习 $H(x)$。这样可以缓解深层网络训练中的退化问题和梯度消失问题，使梯度可以直接通过短路路径传递到前层，从而成功训练上百层甚至更深的网络。
\end{enumerate}

\section{Transformer与attention}
\begin{enumerate}
\item \textbf{Transformer 的自注意力机制原理？}\\
自注意力首先将输入序列映射为查询（Q）、键（K）和值（V），对每个位置 $i$，计算与所有位置 $j$ 的相关性 $a_{ij} \propto \exp(Q_i K_j^\top / \sqrt{d_k})$，再用这些权重对 $V_j$ 加权求和得到输出。这样每个位置都可以直接“看见”序列中所有位置，实现全局依赖建模，并且可并行计算。

\item \textbf{Q、K、V 分别是什么含义？}\\
Q（Query）表示当前要“查询”的位置的表示；K是所有位置的“键”，决定被关注程度；V是对应位置信息的具体内容。注意力权重由 Q 与 K 的相似度决定，再用这些权重对 V 进行加权求和得到新的表示。

\item \textbf{multi-head attention 为什么有效？}\\
多头注意力通过将 Q、K、V 在通道维上分块为多个子空间，每个头在不同子空间独立学习注意力模式。不同头可以关注不同位置、不同子模式或不同关系，提高模型的表示能力。最后将各头输出拼接再线性变换，综合多种视角信息。

\item \textbf{Attention 的优点相较于 RNN？}\\
RNN 需要顺序计算，难以并行，对长距离依赖建模困难且有梯度消失问题。自注意力可以对整个序列做全局依赖建模，每一层的所有位置可并行计算，训练效率高；对长距离依赖更友好，且注意力权重提供可解释性。

\item \textbf{位置编码的作用？}\\
自注意力不包含任何位置信息，如果不加以处理，序列的顺序会完全被忽略。位置编码为每个位置添加一个向量，使模型能够区分不同位置并利用相对/绝对位置信息。常见方案有固定的正余弦位置编码和可学习的位置嵌入。
\end{enumerate}

\section{蒸馏、迁移、对比学习}
\begin{enumerate}
\item \textbf{什么是模型蒸馏？}\\
模型蒸馏通过训练一个小模型去模仿大模型的软输出，如通过软化的 softmax 概率分布。学生模型在学习过程中不仅接收硬标签，还从老师的概率分布中获取类别间“暗知识”，从而在体积更小的情况下保留较多性能。

\item \textbf{什么是迁移学习？}\\
迁移学习通过在源任务/源数据上预训练模型，再将其参数迁移到目标任务中，通过微调少量层或全部网络来适应新任务。这样可以在目标数据较少的情况下，利用源任务中学到的通用特征（如图像边缘、纹理、语言模式等），提升性能并减少训练时间。

\item \textbf{什么是对比学习？InfoNCE 是什么？}\\
对比学习通过构造正样本对（语义相同或相似的样本）和负样本对，使模型拉近正对的表示、推远负对的表示，从而在无监督/自监督情况下学习表示。InfoNCE 是一种常用的对比损失，将问题转化为一个多分类任务，在一组候选中区分正样本，理论上与最大化互信息相关。
\end{enumerate}

\section{评估指标与可解释性}
\begin{enumerate}
\item \textbf{什么是 AUC？如何解释？}\\
AUC（Area Under ROC Curve）是 ROC 曲线下的面积，等价于随机抽取一个正样本和一个负样本时，模型将正样本得分排在负样本前的概率。AUC 不依赖具体阈值，适合评估排序能力，尤其在类别不平衡问题中比精度更有意义。

\item \textbf{ROC 与 PR 曲线的区别？}\\
ROC 曲线以 TPR（召回）为纵轴、FPR 为横轴；PR 曲线以精确率为纵轴、召回为横轴。类别极度不平衡时，ROC 可能看起来不错（FPR 很低），但 PR 曲线能更真实地反映模型对正类的识别性能。因此在正负样本极不平衡时更推荐使用 PR 曲线和 AUPR。

\item \textbf{什么是混淆矩阵？}\\
混淆矩阵在二分类中通常包括：TP、FN、FP、TN，表示预测结果与真实标签的组合计数。基于混淆矩阵可以定义精确率、召回率、F1 等各种评估指标，也可扩展到多分类的情形。

\item \textbf{什么是 F1-score？为什么用？}\\
F1 是精确率和召回率的调和平均：
\[
F1 = \frac{2 \cdot P \cdot R}{P+R}.
\]
当精确率和召回率存在权衡时，F1 提供一个综合指标，只有当两者都高时 F1 才会高，因此特别适用于类别不平衡的分类任务。

\item \textbf{如何判断特征重要性？}\\
线性模型中可以直接看系数大小（在适当标准化之后）；树模型中常通过基于分裂带来的信息增益/不纯度下降计算特征重要性；使用因果推断方法。
\end{enumerate}

\end{document}
